% !PRINTED 435
\weberpage{435}
% Unnumbered chapter-opening leaf; scan sequence fixes this leaf as i0466.

\begin{center}
\gesperrt{Achtzehnter Abschnitt.}

\gesperrt{Galoissche Gruppe der Klassengleichung.}
\end{center}

\begin{center}\rule{0.18\linewidth}{0.4pt}\end{center}

\begin{center}
\textbf{\S\,120. \gesperrt{Relationen zwischen den Klasseninvarianten derselben
Diskriminante.}}
\end{center}

Es ist für die Folge eine Bezeichnung zweckmäßig, durch
die die Abhängigkeit der Klasseninvariante von der \gesperrt{Formenklasse}
einfacher ausgedrückt wird. Wir setzen daher, wenn $k$ eine beliebige
 primitive Klasse der Diskriminante $-m$ bedeutet, die
nach § 118, III. durch $(a, b, c p^{\varepsilon})$ repräsentiert wird:
\[
(k) = j\left(\frac{-b + \sqrt{-m}}{2a}\right) = j(\omega).
\]
Wenn eine oder mehrere der $(k)$ unter einem Funktionszeichen
 auftreten, so werden wir die Klammern auch weglassen,
also z. B. $f(k, k', \ldots)$ für $f[(k), (k'), \ldots]$ schreiben.

Ist also
\weq{1}{(k) = j(\omega),}
und $l$ eine durch $(p, b, a c p^{\varepsilon-1})$ repräsentierte Klasse, so ist nach
§ 118, III.:
\weq{2}{(lk) = j\left(\frac{\omega}{p}\right),}
und es genügen also die beiden Größen
\[
u = (k), \quad v = (lk)
\]
der Invariantengleichung
\weq{3}{F_p(v, u) = 0.}
Die Größe $v = (lk)$ ist aber außerdem eine Wurzel der
Klassengleichung
\weq{4}{H_m(v) = 0,}
und es ist zunächst festzustellen, wieviele Wurzeln die Gleichung
(3) mit (4) gemein hat.
